% !TEX root = ../main.tex

\section{Results}
\label{ch:results}

% Experimental results and analysis. Theoretical results, if any. Meaning of the results. Own experiments, if any.

The authors validate their approach by conducting experiments on various tasks, demonstrating the effectiveness and interpretability of Transformer Programs.

\subsection{In-Context Learning}
\label{subsec:incontext-results-analysis}
The authors initially demonstrate Transformer Programs on a simple in-context learning task\cite{brown2020language}, where the model processes a sequence of alternating letters and numbers and outputs the corresponding number if the last letter has appeared before or "unk" if it has not. The model consists of two layers with one attention head per layer, and its final embedding includes four variables: tokens, positions, and an output variable for each attention head. It achieves perfect accuracy on the test set. The residual streams and attention heads are shown in \cref{fig:incontext} and \cref{fig:incontext-attention}, with the corresponding Python code in \cref{fig:incontext-python}.

\begin{figure}[!ht]
  \centering
  \includegraphics[width=1\textwidth]{Results/incontext-python.png}
  \caption{Figure shows the equivalent Python code for the attention heads (\cref{fig:incontext-attention}) of the model trained on the in-context learning task.}
  \label{fig:incontext-python}
\end{figure}

As shown in \cref{fig:incontext-attention,fig:incontext-python}, the first attention head reads position variable for both key and query At each position, it attends to the key at the previous position and writes the value of the tokens variable as the value. The second attention head uses tokens as the query and the first head's output as the key, matching a letter in the query with a corresponding key that has the same letter followed by a number.

\subsection{RASP Tasks}

\begin{table}[!ht]
  \small
  \centering
  \begin{tabularx}{\textwidth}{l  p{2.95cm}  p{3.7cm}  c  c  c  c  S  S[table-format=2.2]}
    \toprule
    \textbf{Dataset} & \textbf{Description} & \textbf{Example}  & \textbf{k} & \textbf{L} & \textbf{H} & \textbf{M} & \textbf{Acc.} & \textbf{Acc.*}\\
    \midrule
    Reverse & Reverse a string. &  $reverse("abbc") = "cbba"$ & 8 & 2 & - & - & 99.79 & 98.04  \\
    Histogram & For each token, the number of occurrences of that letter in the sequence. & $hist("abbc") = "1221"$ & 8 & 1 & 4 & 2 & 100.00 & 44.24\\
    Double hist. & For each token, the number of unique tokens with the same histogram value. & $hist2("abbc") = "2112"$ & 8 & 3 & 4 & 2 & 98.40 & 65.35 \\
    Sort & Sort the input in lexicographical order. & $sort("cbab") = "abbc"$ & 8 & 3 & 8 & 4 & 99.83 &  99.26\\
    Most-Freq & The unique input tokens in order of frequency, using position to break ties. & $most\_freq("abbc") = "bac"$ & 8 & 3 & 8 & 4 & 75.69 & 72.39\\
    % Dyck-1 & For each position i, is the input up until i a valid string in Dyck-1 $(T)$; a valid prefix $(P)$; or invalid $(F)$. & $dyck1("()())")= "PTPTF"$ & 16 & 3 & 8 & 2 & 99.30 \\
    % Dyck-2 & The same as above, but in Dyck-2. & $dyck2("({})(}") = "PPPTPF"$ & 16 & 3 & 4 & 4 & 99.09 \\
    \bottomrule
  \end{tabularx}
  \caption{Results of RASP tasks introduced by \textcite{weiss2021thinking} on the Transformer Programs including the variable cardinality $(k)$, number of layers $(L)$, attention heads $(H)$, and MLPs $(M)$ in the best model. $(Acc.)$ denotes reported accuracy, while $(Acc.*)$ is the accuracy that I got by training the model for 150 epochs with identical hyperparameters.}
  \label{tab:rasp-task-results}
\end{table}

These are algorithmic tasks introduced by \textcite{weiss2021thinking}. As shown in \cref{tab:rasp-task-results}, the learned programs achieve over 99\% accuracy on three out of five tasks, demonstrating their capability to handle diverse problems. My reimplementation using the same hyperparameters for 150 epochs yielded comparable results, except for the Histogram and Double Histogram tasks. Standard Transformers generally outperform Transformer Programs, particularly on tasks with longer sequences and larger vocabularies (\cref{fig:comp-with-std}). Both models exhibit performance declines as task complexity increases, but Transformer Programs face sharper drops. For the $Reverse$ task with a larger vocabulary, manually initializing attention and MLP weights enabled the model to learn successfully, highlighting the need for further optimization to match standard Transformers on complex tasks.

\begin{figure}[!ht]
    \centering
    \includegraphics[width=1\textwidth]{Results/comp-with-std.png}
    \caption{This shows the comparison of Transformer Programs and standard Transformers on RASP tasks with varying vocabulary size$(|V|)$ and maximum sequence length$(N)$. The performance declines for both models with longer sequences and larger vocabularies, but Transformer Programs show a steeper drop, with accuracy decreasing by over 50\% in the most challenging cases}
    \label{fig:comp-with-std}
\end{figure}


\subsection{NLP Tasks: Named Entity Recognition and Text Classification}

The authors evaluated Transformer Programs on the CoNLL-2003 Named Entity Recognition task \cite{sang2003introduction}, where words are labeled as one of four entity types (PER, ORG, LOC, MISC) or as non-entities (O). The model, trained with categorical attention heads, achieved results comparable to standard Transformers and baselines, as shown in \cref{tab:ner-results}. The analysis highlights the model's ability to leverage contextual information and combine word-level features with cues from neighboring words, such as prepositions, demonstrating its interpretability.
\begin{table}[!ht]
    \small
    \centering
    \begin{tabular}{l c c c S S S[table-format=2.1]}
      \toprule
      \textbf{Model} & \textbf{D} & \textbf{L}& \textbf{H}& \textit{Precision} & \textit{Recall} & \textit{F1} \\
      \midrule
      Unigram baseline & - & - & - & 82.8 & 55.1 & 66.2 \\
      Standard Transformer & 128 & 3 & 4 & 80.4 & 62.7 & 70.5  \\
      Transformer Programs & $32 X 4$ & 2 & 8 & 81.0 & 71.8 & 76.1 \\
      \bottomrule
    \end{tabular}
    \caption{Results and comparisons of baseline, Transformer Programs and standard Transformers on CoNLL-2003 Named Entity Recognition task. For each model, embedding dimension $(D)$, number of layers $(L)$, number of heads $(H)$ of the best-performing model,
    and precision, recall and F1 are repoted. }
    \label{tab:ner-results}
  \end{table}

  The authors also tested Transformer Programs on text classification datasets, including TREC (question classification) \cite{voorhees2000building}, MR (movie review sentiment analysis) \cite{pang2005seeing}, Subj (subjectivity classification) \cite{pang2004sentimental}, and AG News (news topic classification) \cite{zhang2015character}. As shown in \cref{tab:tc-results}, Transformer Programs achieved competitive accuracy compared to standard Transformers, demonstrating their effectiveness in diverse text classification tasks.

\begin{table}[!ht]
    \small
    \centering
    \begin{tabular}{l S S S S[table-format=2.1]}
      \toprule
      \textbf{Model} & \textbf{TREC} & \textbf{MR} & \textbf{Subj} & \textbf{AG} \\
       \midrule
        Bag of words & 74.8 & 77.8 & 92.6 & 89.6 \\
        Standard Transformer & 88.6 & 77.2 & 92.3 & 91.7 \\
        Transformer Program & 85.6 & 77.9 & 93.0 & 90.8 \\
      \bottomrule
    \end{tabular}
    \caption{Classification accuracy on question classification $(TREC)$, sentiment analysis $(MR)$, subjectivity classification $(Subj)$, and topic classification $(AG news)$. The Bag of words baseline is a multinomial Naive Bayes model trained on unigram features.}
    \label{tab:tc-results}
  \end{table}


Overall, The results for NER and text classification tasks highlight Transformer Programs' ability to learn effective and interpretable solutions for complex NLP problems. With performance comparable to standard Transformers and enhanced interpretability, this approach shows promise for developing more transparent and trustworthy AI systems.

